\documentclass[french]{article}
\usepackage[utf8]{inputenc}
\usepackage[T1]{fontenc}
\usepackage{lmodern}
\usepackage{geometry}
\usepackage{graphicx}
\usepackage{babel}
\usepackage{amsmath}
\usepackage{amsthm}
\usepackage{amssymb}
\usepackage{hyperref}
\usepackage{stmaryrd}
\usepackage{enumitem}
\usepackage{tcolorbox}
\usepackage{dsfont}
\usepackage{multirow}
\usepackage{etoc}
\usepackage{titlesec}
\usepackage{esint}
\usepackage{cancel}
\usepackage{mathdots}

\setlist[itemize]{label=$\bullet$}


\renewcommand{\P}{\mathbb{P}}
\newcommand{\N}{\mathbb{N}}
\newcommand{\Z}{\mathbb{Z}}
\newcommand{\Q}{\mathbb{Q}}
\newcommand{\R}{\mathbb{R}}
\newcommand{\C}{\mathbb{C}}
\newcommand{\K}{\mathbb{K}}
\renewcommand{\L}{\mathbb{L}}
\newcommand{\U}{\mathbb{U}}
\newcommand{\st}{^*}
\newcommand{\tq}{\middle|}
\newcommand{\inv}{^{-1}}
\newcommand{\E}{\mathbb{E}}
\newcommand{\F}{\mathbb{F}}
\newcommand{\V}{\mathbb{V}}
\renewcommand{\ker}{\text{Ker}}
\newcommand{\im}{\text{Im}}
\newcommand{\card}{\text{Card}}
\newcommand{\Tr}{\text{Tr}}

\title{TD Euclidiens\\ Exercice 16}
\date{}
\begin{document}
\maketitle

\section{Énoncé}
Soit $u\in\mathcal{L}(E)$. Montrer l'équivalence : $$\left(u^2=0\ \ \text{et} \ \ u+u\st\in\mathcal{GL}(E)\right)\iff \ker(u)=\im(u)$$
\section{Solution}
On raisonne par double implication. Montrons un premier lemme. Supposons $\im(u)\subset\ker(u)$ et montrons que $$\ker(u+u\st)=\ker(u)\cap\ker(u\st)$$ L'inclusion réciproque est immédiate. Dans l'autre sens, supposons $u(x)+u\st(x)=0$. Alors $u\st(x)\in\im(u)$ donc $u(u\st(x))=0$. On en déduit que $(x\mid u(u\st(x)))=0$. Puis par propriété de l'adjoint, $\lVert u\st(x)\rVert=0$ donc $u\st(x)=0$ et donc $u(x)=0$, ce qui est bien la propriété recherchée.
\begin{itemize}
	\item Sens direct : la première hypothèse fournit $\im(u)\subset\ker(u)$. La deuxième fournit $\ker(u+u\st)=\{0\}$. Ensuite, prenons $x\in\ker(u)$. Par surjectivité de $u+u\st$, on peut écrire $x=u(y)+u\st(y)$. Or, $u(x)=0$ et $u(x)=u(u\st(y))$ car $u^2=0$. On en déduit que $u\st(y)=0$ donc que $x\in\im(u)$. 
	\item Sens réciproque : l'inclusion $\im(u)\in\ker(u)$ fournit $u^2=0$ et le lemme donne que si $x\in\ker(u+u\st)$, alors il existe $y$ tel que $x=u(y)$. Or, $u\st(x)=0$ donc $u\st(u(y))=0$ donc $u(y)=0$ donc $x=0$ donc $u+u\st$ est injective donc bijective en dimension finie.
\end{itemize}

\end{document}
